\setcounter{dang}{0}
\setcounter{section}{17}
\section{LŨY THỪA VỚI SỐ MŨ THỰC}
\subsection{Lý thuyết cần nhớ}
\subsubsection{Lũy thừa với số mũ nguyên}
	\begin{itemize}
		\item [\iconCH] Lũy thừa với số mũ nguyên dương: Cho $a\in \mathbb{R}, n\in {\mathbb{N}}^{*}$, khi đó: $a^n=\underbrace{a.a.a...a}_{n \text{ thừa số}}$.
		\item [\iconCH] Lũy thừa với số mũ nguyên âm: Cho $a \ne 0, n\in {\mathbb{N}}^{*}$, khi đó: $a^{-n}=\dfrac{1}{a^n}$.
		\item [\iconCH] Với $a \ne 0$, ta quy ước $a^0=1$; \quad $0^0$ và $0^{-n}\, (n \in \mathbb{N}^{*})$ không có nghĩa.
	\end{itemize}

\subsubsection{Căn bậc n}
\begin{enumerate}
	\item[\iconCH] \indamm{Định nghĩa:} Cho số thực $a$ và số nguyên dương $n$ ($n\ge 2$). Số $b$ được gọi là \textit{căn bậc $n$} của số $a$ nếu $b^n=a$.
	\begin{note}
		\begin{itemize}
			\item[\ding{172}] Với $n$ lẻ và $a\in \mathbb{R}$: Có duy nhất một căn bậc $n$ của $a$, ký hiệu là $\sqrt[n]{a}$.
			\item[\ding{173}] Với $n$ chẵn, ta xét ba trường hợp sau:
			\begin{itemize}
				\item $a<0$: Không tồn tại căn bậc $n$ của $a$;
				\item $a=0$: Có một căn bậc $n$ của $a$ là số $0$;
				\item $a>0$: Có hai căn bậc $n$ của $a$ là hai số đối nhau, kí hiệu giá trị dương là $\sqrt[n]{a}$, còn giá trị âm là $-\sqrt[n]{a}$.
			\end{itemize}
		\end{itemize}
	\end{note}
	\item[\iconCH] \indamm{Tính chất:}
	\begin{boxdn}
	\begin{listEX}[2]
		\item![\ding{172}] $\sqrt[n]{a^n}=\heva{&a\text{ nếu } n \text{ lẻ}\\&|a|\text{ nếu } n \text{ chẵn};}$
		\item[\ding{173}] $\sqrt[n]{a}\cdot \sqrt[n]{b}=\sqrt[n]{ab}$;
		\item[\ding{174}] $\dfrac{\sqrt[n]{a}}{\sqrt[n]{b}}=\sqrt[n]{\dfrac{a}{b}}$;
		\item[\ding{175}] $\left(\sqrt[n]{a}\right)^m=\sqrt[n]{a^m}$;
		\item[\ding{176}] $\sqrt[n]{\sqrt[k]{a}}=\sqrt[nk]{a}$.
	\end{listEX}
	\end{boxdn}
	(\textit{Ở mỗi công thức trên, ta giả sử các biểu thức xuất hiện trong đó có nghĩa}).
\end{enumerate}
\subsubsection{Lũy thừa với số mũ hữu tỉ}
	Cho $a>0$ và số hữu tỉ $r=\dfrac{m}{n}$; trong đó $m\in \mathbb{Z},\,n\in \mathbb{N},\,n\ge 2$. Khi đó: $a^r={a}^{\tfrac{m}{n}}=\sqrt[n]{a^m}.$
\subsubsection{Công thức biến đổi lũy thừa cần nhớ}
	Cho cơ số $a,b>0$ và hai số thực $m,n$. Khi đó, ta có:
	\begin{khung4}{Ghi nhớ}
		\begin{listEX}[3]
			\item [\ding{172}] $a^0=1$; $a^1=a$.
			\item [\ding{173}] $a^{-1}=\dfrac{1}{a}$; $a^{-n}=\dfrac{1}{a^n}$.
			\item [\ding{174}] $\sqrt{a}=a^{\tfrac{1}{2}}$; $\sqrt[n]{a^m}={a}^{\tfrac{m}{n}}$.
			\item [\ding{175}] $a^{m+n}=a^m\cdot a^n$.
			\item [\ding{176}] $a^{m-n}=\dfrac{a^m}{a^n}$.
			\item [\ding{177}] $a^{m\cdot n}=\left(a^m\right)^n=\left(a^n\right)^m$.
			\item [\ding{178}] $(ab)^n=a^n \cdot b^n$.
			\item [\ding{179}] $\left(\dfrac{a}{b}\right)^n=\dfrac{a^n}{b^n}$.
			\item [\ding{180}]  $\left(\dfrac{a}{b}\right)^n=\left(\dfrac{b}{a}\right)^{-n}$.
		\end{listEX}
	\end{khung4}
\subsubsection{So sánh hai lũy thừa cùng cơ số}
Cho cơ số $a>0$ và hai số thực $x,y$. Khi đó, ta có:
\begin{listEX}[2]
	\item [\ding{172}] Nếu $a>1$ thì $a^x>a^y\Leftrightarrow x>y$.
	\item [\ding{173}] Nếu $0<a<1$ thì $a^x>a^y\Leftrightarrow x<y$.
\end{listEX}

\subsection{Các dạng toán thường gặp}
\begin{dang}{Tính giá trị biểu thức}
\end{dang}

\begin{vd}%[1K6BH-1] 
	Tính
	\begin{enumEX}[a)]{3}
	\item $ 25^{\tfrac{3}{2}} $.
	\item $ 32^{-\tfrac{2}{5}} $.
	\item $\left(\dfrac{27}{8}\right)^{\tfrac{2}{3}}$.
	\item $ \sqrt[3]{-27} $.
	\item $\sqrt[5]{9} \cdot \sqrt[5]{27}$.
	\item $\dfrac{\sqrt[3]{128}}{\sqrt[3]{2}}$
	\end{enumEX}
	\loigiai{
		\begin{enumerate}[a)]
		\item $ 25^{\tfrac{3}{2}}=\left(5^2\right)^{\tfrac{3}{2}}=5^3=125 $.
		\item $ 32^{-\tfrac{2}{5}}=\left(2^5\right)^{-\tfrac{2}{5}}=2^{-2}=\dfrac{1}{4} $.
		\item $\left(\dfrac{27}{8}\right)^{\tfrac{2}{3}}=\left(\dfrac{3}{2}\right)^{3\cdot \dfrac{2}{3}}=\left(\dfrac{3}{2}\right)^2=\dfrac{9}{4}$.
		\item $ \sqrt[3]{-27}=\sqrt[3]{(-3)^3}=-3$.
		\item $\sqrt[5]{9} \cdot \sqrt[5]{27}=\sqrt[5]{3^2} \cdot \sqrt[5]{3^3}=\sqrt[5]{3^2 \cdot 3^3}=\sqrt[5]{3^5}=3$;
		\item $\dfrac{\sqrt[3]{128}}{\sqrt[3]{2}}=\sqrt[3]{\dfrac{128}{2}}=\sqrt[3]{64}=\sqrt[3]{4^3}=4$.
		\end{enumerate}
	}
\end{vd} \dongcham{12}

\begin{vd}
	Tính
	\begin{tasks}(1)
		\task $A=\left(\sqrt{2}\right)^{86}-2^{43}$.
		\task $B=\left(\dfrac{1}{16}\right)^{-0,75}+\left(\dfrac{1}{8}\right)^{-\tfrac{4}{3}}$.
		\task $C=\left(\dfrac{3}{7}\right)^{-1}-\dfrac{3}{4}\cdot \left(\dfrac{9}{4}\right)^{-1}$.
		\task $D = \dfrac{6^{3+\sqrt{5}}}{2^{2+\sqrt{5}}\cdot 3^{1+\sqrt{5}}}$.
		\task $E=\left(\dfrac{1}{625}\right)^{-\tfrac{1}{4}}+16^{\tfrac{3}{4}}-2^{-2}.64^{\tfrac{1}{3}}$.
		\task $F = 64^\frac{1}{2}\cdot 64^\frac{1}{3}\cdot \sqrt[6]{64}+{4}^{2+2\sqrt[3]{5}}:{16}^{\sqrt[3]{5}}$
	\end{tasks}
\loigiai{
\begin{enumerate}[a)]
	\item $A=\left(\sqrt{2}\right)^{86}-2^{43}=2^{43}-2^{43}=0$.
	\item $B=\left(\dfrac{1}{16}\right)^{-0,75}+\left(\dfrac{1}{8}\right)^{-\tfrac{4}{3}}=\left(2^{-4}\right)^{-0 \cdot 75}+\left(2^{-3}\right)^{-\tfrac{4}{3}}=2^3+2^4=8+16=24$. 
	\item Biến đổi, ta được $C=\dfrac{7}{3}-\dfrac{3}{4}\cdot \dfrac{4}{9}=2$
	\item Ta có $D = \dfrac{6^{3+\sqrt{5}}}{2^{2+\sqrt{5}}\cdot 3^{1+\sqrt{5}}}=\dfrac{2^{3+\sqrt{5}}\cdot 3^{3+\sqrt{5}}}{2^{2+\sqrt{5}}\cdot 3^{1+\sqrt{5}}}=2\cdot 3^2 =18$.
	\item $E =\left(\dfrac{1}{625}\right)^{\tfrac{-1}{4}}+16^{\tfrac{3}{4}}-2^{-2}
	\cdot 64^{\tfrac{1}{3}} = \left(5^{-4}\right)^{\tfrac{-1}{4}}+(2^4)^{\tfrac{3}{4}}-2^{-2}
	\cdot (2^6)^{\tfrac{1}{3}}=5+2^3-1=12$.
	\item $F = 8\cdot 4\cdot 2 = 64 -4^2=48.$
\end{enumerate}}
\end{vd} \dongcham{18}

\begin{vd}
	Biết rằng $3^x=2$. Tính giá trị của biểu thức: 
	\begin{tasks}(1)
		\task $A=3^{2x-1}\cdot{\left(\dfrac{1}{3}\right)}^{2x-1}+9^{x+1}$.
		\task $B=81^x+\sqrt[4]{3}^x \cdot \sqrt[4]{27^x}$
	\end{tasks}
	\loigiai{
		$A=3^{2x-1}\cdot\left(\dfrac{1}{3}\right)^{2x-1}+9^{x+1}=1+9\cdot(3^x)^2=1+9\cdot4=37$.}
\end{vd} \dongcham{11}

\begin{dang}{Rút gọn biểu thức liên quan đến lũy thừa}
\end{dang}

\begin{vd}%[1D4B1-1]
	Viết mỗi biễu thức sau dưới dạng một luỹ thừa ($a>0$)
	\begin{tasks}(3)
		\task $\sqrt[4]{2^{-3}}$;
		\task $\dfrac{1}{\sqrt[5]{2^3}}$;
		\task $(\sqrt[5]{3})^4$;
		\task $\sqrt{a \sqrt[3]{a}}$;
		\task* $\sqrt[3]{a} \cdot \sqrt[4]{a^3}:(\sqrt[6]{a})^5$;
		\task $a^{\tfrac{1}{3}}: a^{-\tfrac{3}{2}} \cdot a^{\tfrac{2}{3}}$.
	\end{tasks}
	\loigiai{
		\begin{listEX}[1]
			\item $\sqrt[4]{2^{-3}}=2^{-\tfrac{3}{4}}$;
			\item $\dfrac{1}{\sqrt[5]{2^3}}=\dfrac{1}{2^{\tfrac{3}{5}}}=2^{-\tfrac{3}{5}}$;
			\item $(\sqrt[5]{3})^4=3^{\tfrac{4}{5}}$;
			\item $\sqrt{a \sqrt[3]{a}}=\left(a \cdot a^{\tfrac{1}{3}}\right)^{\tfrac{1}{2}}=\left(a^{1+\tfrac{1}{3}}\right)^{\tfrac{1}{2}}=\left(a^{\tfrac{4}{3}}\right)^{\tfrac{1}{2}}=a^{\tfrac{4}{3} \cdot \tfrac{1}{2}}=a^{\tfrac{2}{3}}$;
			\item $\sqrt[3]{a} \cdot \sqrt[4]{a^3}:(\sqrt[6]{a})^5=a^{\tfrac{1}{3}} \cdot a^{\tfrac{3}{4}}:\left(a^{\tfrac{1}{6}}\right)^5=a^{\tfrac{1}{3}+\tfrac{3}{4} -\tfrac{5}{6}}=a^{\tfrac{3}{12}}=a^{\tfrac{1}{4}}$;
			\item $a^{\tfrac{1}{3}}: a^{-\tfrac{3}{2}} \cdot a^{-\tfrac{2}{3}}=a^{\tfrac{1}{3}+\tfrac{3}{2} -\tfrac{2}{3}}=a^{\tfrac{7}{6}}$.
		\end{listEX}	
	}
\end{vd} \dongcham{14}

\begin{vd}
	Viết các biểu thức sau về dạng lũy thừa với số mũ hữu tỉ
	\begin{tasks}(1)
		\task $A =a^{\frac{1}{6}} \sqrt[3]{a}$, với $a>0$.
		\task $B=a^{\frac{5}{3}}\colon \sqrt[3]{a}$, với $a>0.$
		\task $C= \dfrac{ \sqrt[3]{ a^2 \sqrt{a} } }{ a^3 }$, với $a>0$.
		\task $D=\dfrac{a^{\sqrt{7}+1}a^{2-\sqrt{7}}}{(a^{\sqrt{2}-2})^{\sqrt{2}+2}}$ với $a>0$.
	\end{tasks}
	\loigiai{
		\begin{enumerate}[a)]
			\item $A =a^{\frac{1}{6}} \sqrt[3]{a}= a^{\frac{1}{6}}.a^{\frac{1}{3}} = a^{\frac{1}{2}}$. 
			\item  $B=a^{\frac{5}{3}}\colon \sqrt[3]{a}=a^{\frac{5}{3}}\cdot a^{\frac{1}{3}}=a^2.$
			\item $C= a^{\alpha} = \dfrac{ \sqrt[3]{  a^2 \sqrt{a} } }{a^3} = \dfrac{a^{ \frac{5}{6} }}{a^3} = a^{ - \frac{13}{6} }$
			\item $D=\dfrac{a^{\sqrt{7}+1}.a^{2-\sqrt{7}}}{(a^{\sqrt{2}-2})^{\sqrt{2}+2}}=\dfrac{a^{\sqrt{7}+1+2-\sqrt{7}}}{a^{(\sqrt{2}-2)(\sqrt{2}+2)}}=\dfrac{a^3}{a^{-2}}=a^5$.
		\end{enumerate}
}
\end{vd} \dongcham{14}

\begin{dang}{So sánh hai lũy thừa}
\end{dang}

\begin{vd}%[1K6BH-3] 
	Cho $a>0$ và $a \ne 1$. Hãy so sánh cơ số $a$ với $1$,  biết rằng
	\begin{enumEX}[a)]{3}
	\item $ a^{\tfrac{3}{4}} > a^{\tfrac{5}{6}}$.
	\item $a^{\tfrac{11}{6}}< a^{\tfrac{15}{8}}$.
	\item $\sqrt[15]{a^7}>\sqrt[5]{a^2}$.
	\end{enumEX}
	\loigiai{
		\begin{enumerate}
			\item Do $ \dfrac{3}{4}<\dfrac{5}{6} $ và $ a^{\tfrac{3}{4}}>a^{\tfrac{5}{6}} $ nên $ a<1 $.
			\item Do $ \dfrac{11}{6}< \dfrac{15}{8} $ và $ a^{\tfrac{11}{6}}<a^{\tfrac{15}{8}} $ nên $ a<1 $.
			\item Ta có $\sqrt[5]{a^2}=\sqrt[15]{a^6}$. Do đó $\sqrt[15]{a^7}>\sqrt[5]{a^2}\Leftrightarrow \sqrt[15]{a^7}>\sqrt[15]{a^6}
			$ mà $7>6$ nên $a>1$.
		\end{enumerate}
	}
\end{vd} \dongcham{7}

\begin{dang}{Vận dụng, thực tiễn}
\end{dang}


\begin{vd}
	Tính giá trị các biểu thức sau:
	\begin{tasks}(1)
		\task $A=(7+4\sqrt{3})^{2020}(4\sqrt{3}-7)^{2019}$.
		\task $B=\dfrac{\left(4+2\sqrt{3} \right)^{2018 }\cdot \left(1-\sqrt{3} \right)^{2017 }}{\left(1+\sqrt{3} \right)^{2019 }}$.
	\end{tasks}
	\loigiai{
		\begin{enumerate}[a)]
			\item Ta có
			\begin{eqnarray*}
				A&=&(7+4\sqrt{3})(7+4\sqrt{3})^{2019}(4\sqrt{3}-7)^{2019}\\
				&=&1(7+4\sqrt{3})\left[(7+4\sqrt{3})(4\sqrt{3}-7)\right]^{2019}\\
				&=&(7+4\sqrt{3})(-1)^{2019}\\
				&=&-7-4\sqrt{3}.
			\end{eqnarray*}
			\item  \begin{eqnarray*}
				B	&= &\dfrac{\left(4+2\sqrt{3} \right)^{2018 }\cdot \left(1-\sqrt{3} \right)^{2017 }}{\left(1+\sqrt{3} \right)^{2019 }}   \\
				&=&\dfrac{\left(1+\sqrt{3} \right)^{4036 }\cdot \left(1-\sqrt{3} \right)^{2017 }}{\left(1+\sqrt{3} \right)^{2019 }}\\
				&= & \left(1+\sqrt{3} \right)^{2017 }\cdot \left(1-\sqrt{3} \right)^{2017 }\\
				&=&\left[\left(1+\sqrt{3} \right)\cdot \left(1-\sqrt{3} \right) \right]^{2017 }\\
				&=& (-2)^{2017 }=-2^{2017 } .
			\end{eqnarray*}
	\end{enumerate}}
\end{vd} \dongcham{10}

\begin{vd}
	Cho $4^x + 4^{-x} = 14$. Tính giá trị của biểu thức $M = \dfrac{2 + 2^x + 2^{-x}}{7 - 2^x - 2^{-x}}$.
	\loigiai{
		Ta có $4^x + 4^{-x} = 14 \Leftrightarrow \left(2^x + 2^{-x}\right)^2 - 2\cdot 2^x\cdot 2^{-x} = 14 \Leftrightarrow \left(2^x + 2^{-x}\right)^2 = 16 \Rightarrow 2^x + 2^{-x} = 4$.\\
		Vậy $M = \dfrac{2 + 2^x + 2^{-x}}{7 - 2^x - 2^{-x}} = \dfrac{2 + 4}{7 -4} = 2$.
	}
\end{vd} \dongcham{9}

\begin{vd}
	Rút gọn các biểu thức sau (với $a,b>0$ và $a,b \ne 1$):
	\begin{tasks}(1)
		\task $A=\dfrac{a^{\frac{1}{3}}\sqrt{b}+b^{\frac{1}{3}}\sqrt{a}}{\sqrt[6]{a}+\sqrt[6]{b}}$.
		\task $B=\dfrac{a-3a^{\tfrac{1}{3}}+2}{\sqrt[3]{a}-1}+\dfrac{\sqrt{a}-a^{\tfrac{5}{6}}+\sqrt[6]{a}}{\sqrt[6]{a}}$.
	\end{tasks}
	\loigiai{
		\begin{enumerate}[a)]
			\item Ta có $A=\dfrac{a^{\frac{1}{3}}\sqrt{b}+b^{\frac{1}{3}}\sqrt{a}}{\sqrt[6]{a}+\sqrt[6]{b}}=\dfrac{a^{\frac{1}{3}}b^{\frac{1}{3}}\left(a^{\frac{1}{6}}+b^{\frac{1}{6}}\right)}{a^{\frac{1}{6}}+a^{\frac{1}{6}}}=a^{\frac{1}{3}}b^{\frac{1}{3}}=\sqrt[3]{ab}$.
			\item Đặt $t=\sqrt[6]{a}$. Khi đó,
			$$
			B= \dfrac{t^6-3t^2+2}{t^2-1}+\dfrac{t^3-t^5+t}{t}= t^4+t^2-2+t^2-t^4+1= 2t^2-1.$$
			Vậy $B=2\sqrt[3]{a}-1$.
		\end{enumerate}
		
	}
\end{vd} \dongcham{20}

\begin{vd}%[1D4G1-2]
	Cho $x>0, y>0$ và số thực $a$ thoả
	$
	a=\sqrt{x^2+\sqrt[3]{x^4 y^2}}+\sqrt{y^2+\sqrt[3]{x^2 y^4}}.$
	Chứng minh rằng $a^{\frac{2}{3}}=x^{\frac{2}{3}}+y^{\frac{2}{3}}$.
	\loigiai{
		Đặt $\heva{&b=\sqrt[3]{x^2}\\&c=\sqrt[3]{y^2}}\Rightarrow \heva{&x^3=b^3\\&y^2=c^3}$, $(b\ge0, \,c\ge0)$.\\
		Khi đó
		\allowdisplaybreaks
		\begin{eqnarray*}
			a=\sqrt{x^2+\sqrt[3]{x^4 y^2}}+\sqrt{y^2+\sqrt[3]{x^2 y^4}} &\Leftrightarrow& a=\sqrt{b^3+\sqrt[3]{b^6c^3}}+\sqrt{c^3+\sqrt[3]{b^3c^6}}\\
			&\Leftrightarrow& \sqrt{b^3+b^2c}+\sqrt{c^3+bc^2}\\
			&\Leftrightarrow& a=b\sqrt{b+c}+c\sqrt{c+b}\\
			&\Leftrightarrow& a=(b+c)\sqrt{b+c}\\
			&\Leftrightarrow& a=\sqrt{(b+c)^3}\\
			&\Leftrightarrow& \sqrt[3]{a^2}=b+c\\
			&\Leftrightarrow& a^{\frac{2}{3}}=x^{\frac{2}{3}}+y^{\frac{2}{3}}.
		\end{eqnarray*}
	}
\end{vd} \dongcham{16}

\begin{vd}%[1D4B1-1]
	Cường độ ánh sáng tại độ sâu $h$ (m) dưới một mặt hồ được tính bằng công thức $I_h=I_0\left(\dfrac{1}{2}\right)^{\tfrac{h}{4}}$, trong đó $I_0$ là cường độ ánh sáng tại mặt hồ đó.
	\begin{enumerate}
		\item Cường độ ánh sáng tại độ sâu $1$ m bằng bao nhiêu phần trăm so với cường độ ánh sáng tại mặt hồ?
		\item Cường độ ánh sáng tại độ sâu $3$ m gấp bao nhiêu lần cường độ ánh sáng tại độ sâu $6$ m?
	\end{enumerate}
	\loigiai{
		\begin{enumerate}
			\item $\dfrac{I_1}{I_0}=\left(\dfrac{1}{2}\right)^{\tfrac{1}{4}} \approx 0{,}84=84 \%$.
			\item $\dfrac{I_3}{I_6}=\left(\dfrac{1}{2}\right)^{\tfrac{3}{4}-\tfrac{6}{4}}=\left(\dfrac{1}{2}\right)^{-\tfrac{3}{4}} \approx 1{,}68$ (lần).
		\end{enumerate}	
	}
\end{vd} \dongcham{8}

\begin{vd}
	Để dự báo dân số của một quốc gia, người ta sử dụng công thức $S=A(2,71)^{nr}$; trong đó $A$ là dân số của năm lấy làm mốc tính, $S$ là dân số sau $n$ năm, $r$ là tỉ lệ tăng dân số hàng năm. Năm $2017$, dân số Việt Nam là $93~671~600$ người (Tổng cục Thống kê, Niên giám thống kê $2017$, Nhà xuất bản Thống kê, Tr.79). Giả sử tỉ lệ tăng dân số hàng năm không đổi là $0{,}81\%$, dự báo dân số Việt Nam năm $2035$ là bao nhiêu người (kết quả làm tròn đến chữ số hàng trăm)?
	\loigiai{
		Ta có $A=93~671~600$, $r=0{,}81\%$, $n=2035-2017=18$.\\
		Do đó $S=Ae^{nr}=93~671~600\cdot (2,71)^{18\cdot 0{,}81\%}\approx 108~374~700$ (người).\\
		Vậy năm $2035$, nước Việt Nam có khoảng $108~374~700$ người.
	}
\end{vd} \dongcham{7}

\begin{vd}%[1D4T1-1]
	Một chất phóng xạ có chu kì bán rã là $25$ năm, tức là cứ sau $25$ năm, khối lượng của chất phóng xạ đó giảm đi một nửa. Giả sử lúc đầu có $10 \mathrm{~g}$ chất phóng xạ đó. Viết công thức tính khối lượng của chất đó còn lại sau $t$ năm và tính khối lượng của chất đó còn lại sau $120$ năm (làm tròn kết quả đến hàng phần nghìn theo đơn vị gam).
	\loigiai{
		Khối lượng của chất phóng xạ còn lại sau thời gian $t$ (năm) là $$m=m_0\cdot 2^{-\frac{t}{T}},$$ trong đó $m_0$ là khối lượng chất phóng xạ ban đầu, $T$ là chu kì bán rã.\\
		Khối lượng của chất phóng xã còn lại sau $120$ năm là $$m=10\cdot 2^{-\frac{120}{25}}\approx0{,}359\, (\mathrm{g}).$$
	}
\end{vd} \dongcham{8}

\subsection{Bài tập tự luyện}

\begin{bt}%[1D4B1-1]
	Tính giá trị của các biểu thức sau:
	\begin{enumEX}[a)]{3}
		\item  $\sqrt[4]{125} \cdot \sqrt[4]{5}$;
		\item  $\dfrac{\sqrt[4]{243}}{\sqrt[4]{3}}$;
		\item  $\dfrac{\sqrt[3]{3}}{\sqrt[3]{24}}$;
		\item  $\sqrt{\sqrt[3]{64}}$;
		\item $\sqrt[4]{3 \sqrt[3]{3}}$;
		\item  $(-\sqrt[6]{4})^3$.
	\end{enumEX}
	\loigiai{
		\begin{listEX}[1]
			\item $\sqrt[4]{125} \cdot \sqrt[4]{5}=\sqrt[4]{5^3 \cdot 5}=\sqrt[4]{5^4}=5$;
			\item $\dfrac{\sqrt[4]{243}}{\sqrt[4]{3}}=\sqrt[4]{\dfrac{243}{3}}=\sqrt[4]{81}=\sqrt[4]{3^4}=3$;
			\item $\dfrac{\sqrt[3]{3}}{\sqrt[3]{24}}=\sqrt[3]{\dfrac{3}{24}}=\sqrt[3]{\dfrac{1}{8}}=\dfrac{1}{\sqrt[3]{2^3}}=\dfrac{1}{2}$;
			\item $\sqrt{\sqrt[3]{64}}=\sqrt[2\cdot3]{2^6}=\sqrt[6]{2^6}=2$;
			\item $\sqrt[4]{3 \sqrt[3]{3}}=\sqrt[4]{\sqrt[3]{3^3 \cdot 3}}=\sqrt[3]{\sqrt[4]{3^4}}=\sqrt[3]{3}$
			\item $(-\sqrt[6]{4})^3=-\sqrt[6]{4^3}=-\sqrt[6]{2^{2\cdot3}}=-\sqrt[6]{2^6}=-2$.
		\end{listEX}	
	}
\end{bt} \dongcham{12}

\begin{bt}
	Tính giá trị của biểu thức
	\begin{tasks}(1)
		\task $A = 4^4\cdot 8^{11}\cdot 2^{2017}.$
		\task $B=3^{1-\sqrt {2}}\cdot 3^{2+\sqrt{2}}\cdot 9^{\tfrac{1}{2}}$
		\task $C=2018^{\sin^2\alpha} \cdot 2018^{\cos^2\alpha}$
		\task $D = \dfrac{6^{2+\sqrt{5}}}{2^{2+\sqrt{5}}\cdot 3^{1+\sqrt{5}}}$.
	\end{tasks}
	\loigiai{
		\begin{enumerate}[a)]
		\item Ta có $A = 4^4\cdot 8^{11}\cdot 2^{2017}=2^8\cdot 2^{33}\cdot 2^{2017}=2^{2058}$
		\item Ta có $B=3^{1-\sqrt {2}}\cdot 3^{2+\sqrt{2}}\cdot 9^{\tfrac{1}{2}}=3^{1-\sqrt {2}+2+\sqrt {2}+1}=3^4=81$.
		\item Ta có $C=2018^{\sin^2\alpha} \cdot 2018^{\cos^2\alpha}=2018^{\sin^2\alpha +\cos^2\alpha}=2018$.
		\item Ta có $D = \dfrac{6^{2+\sqrt{5}}}{2^{2+\sqrt{5}}\cdot 3^{1+\sqrt{5}}}=\dfrac{3^{2+\sqrt{5}}}{3^{1+\sqrt{5}}}=3$.
\end{enumerate}}
\end{bt} \dongcham{14}

\begin{bt}
	Viết các biểu thức sau về dạng lũy thừa của số mũ hữu tỉ
	\begin{tasks}(1)
		\task $ A=a^{\frac{8}{3}}:\sqrt[3]{a^4}$
		\task $B=a^{\tfrac{1}{3}}\cdot \sqrt[4]{a}$
		\task $C=\dfrac{a^{2}a^{\frac{5}{2}}\sqrt[3]{a^{4}}}{\sqrt[6]{a^{5}}}$ ($a>0$)
		\task $D=\dfrac{7^{5+2020\sqrt{5}} \cdot 7^{5-2020\sqrt{5}}}{\left(7^{\sqrt{3}-2}\right)^{\sqrt{3}+2}}$.
	\end{tasks}
	\loigiai{
		\begin{enumerate}[a)]
			\item $A= a^{\frac{8}{3}}:\sqrt[3]{a^4}=a^{\frac{8}{3}}: a^{\frac{4}{3}}=a^{\frac{8}{3}-\frac{4}{3}}=a^{\frac{4}{3}}$
			\item $B=a^{\tfrac{1}{3}}\cdot a^{\tfrac{1}{4}}=a^{\tfrac{7}{12}}$.
			\item Ta có $C=\dfrac{a^{2}a^{\frac{5}{2}}\sqrt[3]{a^{4}}}{\sqrt[6]{a^{5}}}=\dfrac{a^{2+\frac{5}{2}+\frac{4}{3}}}{a^{\frac{5}{6}}}=\dfrac{a^{\frac{35}{6}}}{a^{\frac{5}{6}}}=a^{5}$.
			\item $D=\dfrac{7^{5+2020\sqrt{5}} \cdot 7^{5-2020\sqrt{5}}}{\left(7^{\sqrt{3}-2}\right)^{\sqrt{3}+2}} = \dfrac{7^{10}}{7^{-1}}=7^{11}$.
	\end{enumerate}}
\end{bt} \dongcham{14}

\begin{bt}%[1K6BH-2] 
	Cho $ a $ là số thực dương. Rút gọn các biểu thức sau:
	\begin{enumEX}[a)]{2}
		\item $\left(a^{\sqrt{6}}\right)^{\sqrt{24}}$.
		\item $ a^{\sqrt{2}}\left(\dfrac{1}{a}\right)^{\sqrt{2}-1}$.
		\item $ a^{-\sqrt{3}}: a^{\left(\sqrt{3}-1\right)^2} $.
		\item $ \sqrt[3]{a}\cdot\sqrt[4]{a}\cdot\sqrt[12]{a^5}$.
	\end{enumEX}
	\loigiai{
		\begin{enumerate}
			\begin{multicols}{2}
				\item $\left(a^{\sqrt{6}}\right)^{\sqrt{24}}=a^{\sqrt{6 \cdot24}}=a^{12}$.
				\item $ a^{\sqrt{2}}\left(\dfrac{1}{a}\right)^{\sqrt{2}-1}=a^{\sqrt{2}}\cdot a^{1-\sqrt{2}} =a $.
				\item $ a^{-\sqrt{3}}\colon a^{\left(\sqrt{3}-1\right)^2}=a^{-\sqrt{3}}\colon a^{4-2\sqrt{3}}=a^{-4+\sqrt{3}} $.
				\item $ \sqrt[3]{a}\cdot\sqrt[4]{a}\cdot \sqrt[12]{a^5}=a^{\tfrac{1}{3}}\cdot a^{\tfrac{1}{4}} \cdot a^{\tfrac{5}{12}}=a$.
			\end{multicols}
		\end{enumerate}
	}
\end{bt} \dongcham{13}

\begin{bt}%[1D4B1-2]
	Biết rằng $5^{2x}=3$. Tính giá trị của biểu thức $\dfrac{5^{3x}+5^{-3x}}{5^x+5^{-x}}$.
	\loigiai{
		\allowdisplaybreaks
		\begin{eqnarray*}
			\dfrac{5^{3 x}+5^{-3 x}}{5^x+5^{-x}}=\dfrac{\left(5^x+5^{-x}\right)\left(5^{2 x}-5^x 5^{-x}+5^{-2 x}\right)}{5^x+5^{-x}}=5^{2 x}-1+5^{-2 x}=3-1+\dfrac{1}{3}=\dfrac{7}{3}.
		\end{eqnarray*}
	}
\end{bt} \dongcham{9}

\begin{bt}%[Đề thi GHK1 Lương Thế Vinh 2020-2021]%[Phạm Tuấn, EX-1-2021]%[2D2B1-2]%
	Cho $9^{x}+9^{-x}=23$. Tính giá trị biểu thức $A=\dfrac{5+3^{x}+3^{-x}}{1-3^{x}-3^{-x}}$.
	\loigiai{
		Ta có $(3^{x}+3^{-x})^2 = 9^x + 9^{-x}+2 =25 \Rightarrow 3^{x}+3^{-x} =5$. \\
		Do đó $A= \dfrac{5+3^{x}+3^{-x}}{1-3^{x}-3^{-x}} = \dfrac{5+5}{1-5} = - \dfrac{5}{2}$.\\
		}
\end{bt} \dongcham{10}

\begin{bt}
	Với $0<a$ và $a \ne 1$, hãy tìm $a$ để
	\begin{enumEX}[a)]{2}
		\item  $(a-1)^{\tfrac{-2}{3}}<(a-1)^{\tfrac{-1}{3}}$
		\item  $(a-1)^{-2}>(a-1)^{\sqrt{2}}$
	\end{enumEX}
	\loigiai{
		\begin{enumerate}[a)]
			\item Ta có $\heva{&-\dfrac{2}{3}<-\dfrac{1}{3}\\&(a-1)^{-\tfrac{2}{3}}<(a-1)^{-\tfrac{1}{3}}}\Rightarrow a-1>1\Rightarrow a>2$
			\item Vì $\heva{& -2<\sqrt{2} \\ & (a-1)^{-2}>(a-1)^{\sqrt{2}}}$ nên $0<a-1<1\Leftrightarrow 1<a<2$.
	\end{enumerate}}
\end{bt} \dongcham{12}

\begin{bt}%[1D4K1-3]
	Xác định các giá trị cùa số thụrc $a$ thoả mãn
	\begin{enumEX}[a)]{3}
		\item $a^{\frac{1}{2}}>a^{\sqrt{3}}$;
	\item $a^{-\frac{3}{2}}<a^{\frac{2}{3}}$;
	\item $(\sqrt{2})^a>(\sqrt{3})^a$.
	\end{enumEX}
	\loigiai{
		\begin{enumerate}
			\item Ta có $\dfrac{1}{2}<\sqrt{3}$ nên $a^{\frac{1}{2}}>a^{\sqrt{3}}$ suy ra $0<a<1$.
			\item Ta có $-\dfrac{3}{2}<\dfrac{2}{3}$ nên $a^{-\frac{3}{2}}<a^{\frac{2}{3}}$ suy ra $a>1$.
			\item Ta có $\sqrt{2}<\sqrt{3}$ nên $(\sqrt{2})^a>(\sqrt{3})^a$ suy ra $a<0$.
		\end{enumerate}
	}
\end{bt} \dongcham{8}

\begin{bt}%[1K6TH-1]
	Nếu một khoản tiền gốc $P$ được gửi ngân hàng với lãi suất hằng năm $r$ ($r$ được biểu thị dưới dạng số thập phân), được tính lãi $n$ lần trong một năm, thì tổng số tiền $A$ nhận được (cả vốn lẫn lãi) sau $N$ kì gửi cho bởi công thức sau:
	$$
	A=P\left(1+\dfrac{r}{n}\right)^N .
	$$
	Hỏi nếu bác An gửi tiết kiệm số tiền $120$ triệu đồng theo kì hạn $6$ tháng với lãi suất không đổi là $5 \%$ một năm, thì số tiền thu được (cả vốn lẫn lãi) của bác An sau $2$ năm là bao nhiêu?
	\loigiai{
		Ta có $2$ năm là $24$ tháng ứng với $N=4$ kì hạn.\\
		Do kì hạn là $6$ tháng nên mỗi năm được tính lãi $n=2$ lần.\\
		Vậy số tiền cả vốn lẫn lãi bác An nhận được sau $2$ năm là $A=120\left(1+\dfrac{0{,}05}{2}\right)^4\approx 132{,}457$ triệu đồng.
	}
\end{bt} \dongcham{9}
\begin{bt}
	Định luật thứ ba của Kepler nói rằng bình phương của chu kì quỹ đạo $p$ (tính bằng năm Trái Đất) của một hành tỉnh chuyển động xung quanh Mặt Trời (theo quỹ đạo là một đường elip với Mặt Trời nằm ở một tiêu điểm) bẳng lập phương của bán trục lớn d (tính bằng đơn vị thiên văn $AU$).
	\begin{enumerate}
		\item  Tinh $p$ theo $d$.
		\item  Nếu Sao Thỗ có chu kì quỹ đạo là 29,46 năm Trái Đất, hãy tính bán trục lớn quỹ đạo của Sao Thổ đến Mặt Trời (\textit{kết quả tính theo đơn vị thiên văn và làm tròn đến hàng phần trăm}).
	\end{enumerate}
	\loigiai{
		\begin{enumerate}
			\item  Theo định luật thứ ba của Kepler, ta có:\\
			$$
			p^2=d^3 \text {hay} p=\sqrt{d^3}.
			$$
			
			\item  Thay $p=29,46$ vào công thức $p=\sqrt{d^3}$, ta được $d \approx 9,54 AU$.
		\end{enumerate}
	}
\end{bt} \dongcham{11}

\begin{bt}%[1K6BH-2] 
	Khoảng cách từ một hành tinh đến Mặt Trời có thể xấp xỉ bằng một hàm số của độ dài năm của hành tinh đó. Công thức của hàm số đó là $d=\sqrt[3]{6 t^2}$, trong đó $d$ là khoảng cách từ hành tinh đó đến Mặt Trời (tính bằng triệu dặm) và $t$ là độ dài năm của hành tinh đó (tính bằng số ngày Trái Đất).\\
	\begin{flushright}
		\textit{(Theo Algebra 2, NXB MacGraw-Hill, 2008).}
	\end{flushright}
	\begin{enumerate}
		\item  Nếu độ dài của một năm trên Sao Hoả là 687 ngày Trái Đất thì khoảng cách từ Sao Hoả đến Mặt Trời là bao nhiêu?
		\item  Tính khoảng cách từ Trái Đất đến Mặt Trời (coi một năm trên Trái Đất có 365 ngày).
	\end{enumerate}
	\loigiai{
		\begin{enumerate}
			\item  Thay $t=687$ vào công thức ta được khoảng cách tứ Sao Hoả đến Mặt Trời là
			$$
			d=\sqrt[3]{6 t^2}=\sqrt[3]{6 \cdot 687^2} \approx 141,48 \text {(triệu dặm)}.
			$$
			\item  Thay $t=365$ vào công thức ta được khoảng cách từ Trái Đất đến Mặt Trời là:
		\end{enumerate}
		$$
		d=\sqrt[3]{6 t^2}=\sqrt[3]{6 \cdot 365^2} \approx 92,81 \text {(triệu dặm)}.
		$$
	}
\end{bt} \dongcham{12}

\begin{bt}
	Tại một vùng biển, giả sử cường độ ánh sáng $I$ thay đổi theo độ sâu theo công thức $I=I_0 \cdot 10^{-0,3d}$, trong đó $d$ là độ sâu (tính bằng mét) so với mặt hồ, $I_0$ là cường độ ánh sáng tại mặt hồ.
	\begin{enumerate}
		\item Tại độ sâu $1$ m, cường độ ánh sáng gấp bao nhiêu lần $I_0$?
		\item Cường độ ánh sáng tại độ sâu $2$ m gấp bao nhiêu lần so với tại độ sâu $10$ m? Làm tròn kết quả đến hai chữ số thập phân.
	\end{enumerate}
	\loigiai{
		\begin{enumerate}
			\item Tại độ sâu $d=1$ m $\Rightarrow I=I_0\cdot 10^{-0,3\cdot 1}=I_0\cdot 10^{-0,3}\Rightarrow \dfrac{I}{I_0}=10^{-0,3}=0{,}50$.
			\item Tại độ sâu $d=2$ m $\Rightarrow I=I_0\cdot 10^{-0,3\cdot 2}=I_0\cdot 10^{-0,6}$.\\
			Tại độ sâu $d=10$ m $\Rightarrow I'=I_0\cdot 10^{-0,3\cdot 10}=I_0\cdot 10^{-3}$.\\
			Suy ra $\dfrac{I}{I'}=\dfrac{I_0\cdot 10^{-0,6}}{I_0\cdot 10^{-3}}=\dfrac{10^{-0,6}}{10^{-3}}=10^{-0,6+3}=251{,}19$.
		\end{enumerate}
	}
\end{bt} \dongcham{15}
